\subsection{Qualitative Analysis Results}
\subsubsection{Participant Comment Positivity}
Building on the baseline distribution-and-regression checks reported in Qualitative1, this section develops a more structured profile analysis linking language valence and social orientation to trust-change patterns.

To complement the trust analyses, we estimated the affective tone of participant open responses using lexicon-based sentiment analysis. For each participant, the text fields \emph{AI Interaction Feeling}, \emph{Need Model Understanding}, \emph{Job Screening Feeling}, and \emph{Explanation Comment} were scored with the VADER compound sentiment measure and then averaged to create a participant-level sentiment score. VADER is a rule-based sentiment model designed for short-form natural language that combines lexical valence with heuristics for emphasis, negation, punctuation, and capitalization. In applied work, sentiment measures of this kind are commonly used to summarize reaction valence in online comments, survey responses, social media posts, and user feedback, especially when the goal is to compare whether one group reacts more positively or negatively than another rather than to recover deeper semantic structure.

For visualization, participants were separated into negative and positive sentiment groups using the sign of the mean sentiment score, with neutral responses excluded from the left-panel comparison. Trust change was represented using normalized pre--post change scores for analytical, overall, and emotional trust.

\subsubsection{Participant Comment Individualism and Collectivism}
Because sentiment alone does not capture the social orientation of qualitative responses, we also constructed an individualism--collectivism (IC) text score from the same participant comments. This IC score was computed with weighted lexical dictionaries: first-person singular, autonomy, independence, privacy, and self-oriented terms contributed to the individualistic side, whereas plural pronouns, community, cooperation, collective, mutuality, and shared-responsibility terms contributed to the collectivist side. For each participant, the IC score was defined as a normalized contrast between collectivist and individualistic weights, and the resulting participant-level scores were averaged across the available response fields. Participants were then ordered by this IC score and split into terciles, with the lowest third labeled \emph{individualistic}, the highest third labeled \emph{collectivist}, and the middle third excluded from the direct comparison.

This approach goes beyond conventional sentiment analysis by asking not just whether a response is positive or negative, but what kind of social orientation it reflects. Related work in cultural psychology and computational social science often uses pronoun distributions, autonomy-vs.-interdependence language, moral framing, and semantic association structures to study differences between more individual-focused and more community-focused discourse. Here we use a lightweight lexical approach to test whether the social orientation of participants' comments aligns with the observed trust-change patterns.

\subsubsection{Results}
Figure~\ref{fig:qualitative} summarizes three related analyses. The left panel compares negative and positive sentiment groups on normalized trust change while overlaying the mean trust-change profiles for WEIRD and non-WEIRD participants. The strongest separation appears for emotional trust: participants with negative-sentiment comments showed higher emotional trust change than participants with positive-sentiment comments ($p=.003$), while the analytical and overall contrasts were not statistically significant ($p=.304$ and $p=.356$, respectively). Visually, the WEIRD profile is closer to the negative-sentiment pattern, whereas the non-WEIRD profile is closer to the positive-sentiment pattern.

The middle panel compares participants whose comments fell into the individualistic versus collectivist IC extremes. Here the clearest differences appear in overall and emotional trust change. Collectivist-oriented comments were associated with more positive normalized overall trust change than individualistic comments ($p=.007$), and also with more positive emotional trust change ($p<.001$). The analytical trust contrast was small and non-significant ($p=.360$). As in the sentiment comparison, the WEIRD mean profile visually tracks the individualistic side more closely, while the non-WEIRD mean profile tracks the collectivist side.

The right panel compresses these qualitative patterns into profile-contrast summaries for sentiment, IC orientation, and their combined interaction. Positive values indicate relatively greater similarity to positive-sentiment and individualistic profiles, whereas negative values indicate greater similarity to negative-sentiment and collectivist profiles. Across all three summaries, WEIRD and non-WEIRD participants differ significantly: sentiment contrast ($p<.001$), IC contrast ($p<.001$), and the combined interaction contrast ($p<.001$). The direction of these effects indicates that WEIRD participants align more strongly with the negative/individualistic side of the qualitative profile space, while non-WEIRD participants align more strongly with the positive/collectivist side. Taken together, these results suggest that the qualitative content of participant comments is not merely affective noise, but reflects structured differences in both emotional valence and social orientation that parallel the quantitative trust-change patterns.